% GENERATED FILE. Edit canonical lessons, then run npm run generate:books.
% canonical-source-hash: fnv1a64:cb289e24a5345c2b
% canonical-lessons: GE-C465-sturm, GE-C465-boee, GE-C465-orkan, GE-C465-ueberschwemmen, GE-C465-pfuetze, GE-R465-review-462-465, GE-R465-review-final-465

\chapter{Sturm und Hochwasser}
\label{ch:ge-sturm-und-hochwasser}
\hlchaptermodality{465}

\begin{chapteropening}
\textbf{By the end of this chapter:} \emph{I can talk about stormy weather and floods: the storm, the gust, the hurricane, flooding and the puddle.}

Complete the last lesson of chapter 465: I can talk about stormy weather and floods: the storm, the gust, the hurricane, flooding and the puddle.
\end{chapteropening}

\section[der Sturm]{der Sturm — the storm, the gale}

\label{lesson:GE-C465-sturm}



\begin{quote}
Before the new one: say the German for muggy, then the German for to dry out.
\end{quote}

\subsection*{You'll want to know: der Sturm}
\textbf{der Sturm} — \textquotedblleft{}the storm, the gale\textquotedblright{}.

\begin{quote}
Der Sturm hat einen Baum auf unser Auto geworfen.
\end{quote}

\begin{quote}
\emph{The storm blew a tree down onto our car.}
\end{quote}

Plural: \emph{die Stürme}. \emph{Stürmisch} — stormy.

\begin{cousinweb}
\emph{Sturm} is inherited from Old High German \emph{sturm} 'storm, battle' and Proto-Germanic \emph{sturmaz} (EtymDB); English \emph{storm} is a cousin.
\end{cousinweb}

\begin{grammarlens}[title={What you've built}]
One more word for telling, explaining and describing at B1.
\end{grammarlens}

\subsection*{Your turn}
\begin{itemize}
\raggedright
  \item \emph{Say it:} \emph{der Sturm}
  \item \emph{Say it:} \emph{der Sturm}, once more
  \item \emph{Say it:} a sentence of your own with \emph{der Sturm}
  \item \emph{Recall:} say the German for muggy, then the German for to dry out, then say \emph{der Sturm} again
\end{itemize}

\begin{tcolorbox}[breakable,colback=teal!4,colframe=teal!35!black,title={Before you move on}]
What does der Sturm mean? (\textquotedblleft{}The storm, the gale\textquotedblright{}.) Say it once more.
\end{tcolorbox}
\section[die Böe]{die Böe — the gust (of wind)}

\label{lesson:GE-C465-boee}



\begin{quote}
Before the new one: say the German for to dry out, then the German for the storm.
\end{quote}

\subsection*{You'll want to know: die Böe}
\textbf{die Böe} — \textquotedblleft{}the gust (of wind)\textquotedblright{}.

\begin{quote}
Eine starke Böe nimmt ihr den Regenschirm weg.
\end{quote}

\begin{quote}
\emph{A strong gust takes her umbrella away.}
\end{quote}

Plural: \emph{die Böen}. Also spelled \emph{die Bö}.

\begin{cousinweb}
\emph{Böe} came from Low German and Dutch \emph{bui} 'a sudden gust, a squall', a sailors' word; EtymDB does not list it, and the older origin is not settled.
\end{cousinweb}

\begin{grammarlens}[title={What you've built}]
One more word for telling, explaining and describing at B1.
\end{grammarlens}

\subsection*{Your turn}
\begin{itemize}
\raggedright
  \item \emph{Say it:} \emph{die Böe}
  \item \emph{Say it:} \emph{die Böe}, once more
  \item \emph{Say it:} a sentence of your own with \emph{die Böe}
  \item \emph{Recall:} say the German for to dry out, then the German for the storm, then say \emph{die Böe} again
\end{itemize}

\begin{tcolorbox}[breakable,colback=teal!4,colframe=teal!35!black,title={Before you move on}]
What does die Böe mean? (\textquotedblleft{}The gust (of wind)\textquotedblright{}.) Say it once more.
\end{tcolorbox}
\section[der Orkan]{der Orkan — the hurricane}

\label{lesson:GE-C465-orkan}



\begin{quote}
Before the new one: say the German for the storm, then the German for the gust.
\end{quote}

\subsection*{You'll want to know: der Orkan}
\textbf{der Orkan} — \textquotedblleft{}the hurricane\textquotedblright{}.

\begin{quote}
Der Orkan hat in der Nacht viele Dächer beschädigt.
\end{quote}

\begin{quote}
\emph{The hurricane damaged many roofs in the night.}
\end{quote}

Plural: \emph{die Orkane}. Stronger than a \emph{Sturm}.

\begin{cousinweb}
\emph{Orkan} was borrowed from Dutch \emph{orkaan}, which came from Spanish \emph{huracán}, a word taken from a Caribbean language (EtymDB).
\end{cousinweb}

\begin{grammarlens}[title={What you've built}]
One more word for telling, explaining and describing at B1.
\end{grammarlens}

\subsection*{Your turn}
\begin{itemize}
\raggedright
  \item \emph{Say it:} \emph{der Orkan}
  \item \emph{Say it:} \emph{der Orkan}, once more
  \item \emph{Say it:} a sentence of your own with \emph{der Orkan}
  \item \emph{Recall:} say the German for the storm, then the German for the gust, then say \emph{der Orkan} again
\end{itemize}

\begin{tcolorbox}[breakable,colback=teal!4,colframe=teal!35!black,title={Before you move on}]
What does der Orkan mean? (\textquotedblleft{}The hurricane\textquotedblright{}.) Say it once more.
\end{tcolorbox}
\section[überschwemmen]{überschwemmen — to flood}

\label{lesson:GE-C465-ueberschwemmen}



\begin{quote}
Before the new one: say the German for the gust, then the German for the hurricane.
\end{quote}

\subsection*{You'll want to know: überschwemmen}
\textbf{überschwemmen} — \textquotedblleft{}to flood\textquotedblright{}.

\begin{quote}
Der Fluss überschwemmt die Wiesen am Ufer.
\end{quote}

\begin{quote}
\emph{The river floods the meadows on its banks.}
\end{quote}

Inseparable: \emph{er überschwemmt}, \emph{er hat überschwemmt}.

\begin{cousinweb}
\emph{überschwemmen} is built from \emph{über} 'over' and \emph{schwemmen} 'to float, to wash along', a causative of \emph{schwimmen} 'to swim' (EtymDB lists the verb).
\end{cousinweb}

\begin{grammarlens}[title={What you've built}]
One more word for telling, explaining and describing at B1.
\end{grammarlens}

\subsection*{Your turn}
\begin{itemize}
\raggedright
  \item \emph{Say it:} \emph{überschwemmen}
  \item \emph{Say it:} \emph{überschwemmen}, once more
  \item \emph{Say it:} a sentence of your own with \emph{überschwemmen}
  \item \emph{Recall:} say the German for the gust, then the German for the hurricane, then say \emph{überschwemmen} again
\end{itemize}

\begin{tcolorbox}[breakable,colback=teal!4,colframe=teal!35!black,title={Before you move on}]
What does überschwemmen mean? (\textquotedblleft{}To flood\textquotedblright{}.) Say it once more.
\end{tcolorbox}
\section[die Pfütze]{die Pfütze — the puddle}

\label{lesson:GE-C465-pfuetze}



\begin{quote}
Before the new one: say the German for the hurricane, then the German for to flood.
\end{quote}

\subsection*{You'll want to know: die Pfütze}
\textbf{die Pfütze} — \textquotedblleft{}the puddle\textquotedblright{}.

\begin{quote}
Das Kind springt mit den Stiefeln in jede Pfütze.
\end{quote}

\begin{quote}
\emph{The child jumps into every puddle in his boots.}
\end{quote}

Plural: \emph{die Pfützen}.

\begin{cousinweb}
\emph{Pfütze} is an early borrowing, Old High German \emph{pfuzzi}, from Latin \emph{puteus} 'well, pit' (EtymDB; Lewis \& Short); English \emph{pit} goes back to the same Latin word.
\end{cousinweb}

\begin{grammarlens}[title={What you've built}]
One more word for telling, explaining and describing at B1.
\end{grammarlens}

\subsection*{Your turn}
\begin{itemize}
\raggedright
  \item \emph{Say it:} \emph{die Pfütze}
  \item \emph{Say it:} \emph{die Pfütze}, once more
  \item \emph{Say it:} a sentence of your own with \emph{die Pfütze}
  \item \emph{Recall:} say the German for the hurricane, then the German for to flood, then say \emph{die Pfütze} again
\end{itemize}

\begin{tcolorbox}[breakable,colback=teal!4,colframe=teal!35!black,title={Before you move on}]
What does die Pfütze mean? (\textquotedblleft{}The puddle\textquotedblright{}.) Say it once more.
\end{tcolorbox}
\section[(dialogue)]{Regen, Wind und Sonne}

\label{lesson:GE-R465-review-462-465}



\begin{quote}
Say the German for the severe weather and to come up.
\end{quote}

\subsection*{Your turn}
Say these aloud:

\begin{itemize}
\raggedright
  \item the severe weather, to come up, the hail, to hail, to thunder
  \item the hoarfrost, to freeze, to thaw, the thaw, to shiver
  \item the heatwave, the drought, the humidity, muggy, to dry out
  \item the storm, the gust, the hurricane, to flood, the puddle
\end{itemize}

\begin{tcolorbox}[breakable,colback=teal!4,colframe=teal!35!black,title={Before you move on}]
The severe weather? (\textbf{das Unwetter}.) To come up? (\textbf{aufziehen}.) The hail? (\textbf{der Hagel}.) To hail? (\textbf{hageln}.) To thunder? (\textbf{donnern}.) The hoarfrost? (\textbf{der Raureif}.) To freeze? (\textbf{gefrieren}.) To thaw? (\textbf{tauen}.) The thaw? (\textbf{das Tauwetter}.) To shiver? (\textbf{bibbern}.) The heatwave? (\textbf{die Hitzewelle}.) The drought? (\textbf{die Dürre}.) The humidity? (\textbf{die Luftfeuchtigkeit}.) Muggy? (\textbf{schwül}.) To dry out? (\textbf{austrocknen}.) The storm? (\textbf{der Sturm}.) The gust? (\textbf{die Böe}.) The hurricane? (\textbf{der Orkan}.) To flood? (\textbf{überschwemmen}.) The puddle? (\textbf{die Pfütze}.)
\end{tcolorbox}
\section[(dialogue)]{Sturm und Hochwasser — once more}

\label{lesson:GE-R465-review-final-465}



\begin{quote}
Say the German for the storm and the gust.
\end{quote}

\subsection*{Your turn}
\begin{itemize}
\raggedright
  \item \emph{Say it:} the storm, the gust, the hurricane, to flood, the puddle
\end{itemize}

\begin{tcolorbox}[breakable,colback=teal!4,colframe=teal!35!black,title={Before you move on}]
The storm? (\textbf{der Sturm}.) The gust? (\textbf{die Böe}.) The hurricane? (\textbf{der Orkan}.) To flood? (\textbf{überschwemmen}.) The puddle? (\textbf{die Pfütze}.)
\end{tcolorbox}
